\frontmatter

% ============================================================
% Halaman Sampul
% ============================================================
\begin{titlepage}
  \centering
  \vspace*{1cm}
  {\LARGE\bfseries Kriptografi\par}
  \vspace{0.5cm}
  {\large Buku Ajar Berbasis Outcome-Based Education (OBE)\par}
  \vspace{1cm}
  {\large Mata Kuliah: Kriptografi\par}
  {\large Kode: MAT-XXX\par}
  \vspace{2cm}
  {\large Program Studi Matematika\par}
  {\large Fakultas MIPA / Sains\par}
  {\large 3 SKS, 16 Pertemuan\par}
  \vfill
  {\large 2026\par}
\end{titlepage}

\cleardoublepage

% ============================================================
% Kata Pengantar
% ============================================================
\chapter*{Kata Pengantar}
\addcontentsline{toc}{chapter}{Kata Pengantar}

Puji syukur kehadirat Tuhan Yang Maha Esa atas terselesaikannya buku ajar \textit{Kriptografi} ini. Buku ini disusun dengan pendekatan \textbf{Outcome-Based Education (OBE)}, yang berfokus pada pencapaian kompetensi terukur sesuai dengan Capaian Pembelajaran Lulusan (CPL) dan Capaian Pembelajaran Mata Kuliah (CPMK).

Kriptografi merupakan ilmu yang menggabungkan matematika dan komputasi untuk mengamankan informasi. Mata kuliah ini membekali mahasiswa Program Studi Matematika dengan kemampuan memahami konsep-konsep teoretis kriptografi dan mengimplementasikannya dalam bahasa pemrograman C.

Buku ini dirancang untuk mendukung pembelajaran dengan pendekatan student-centered learning. Setiap bab dilengkapi dengan:
\begin{itemize}
  \item Sub-CPMK yang jelas dan terukur
  \item Materi pokok dengan contoh kode C yang lengkap
  \item Aktivitas pembelajaran yang mendorong eksplorasi mandiri
  \item Latihan dan refleksi untuk penguatan pemahaman
  \item Asesmen untuk mengukur pencapaian kompetensi
  \item Checklist kompetensi untuk self-assessment
\end{itemize}

Kami berharap buku ini dapat menjadi panduan yang efektif dalam perjalanan pembelajaran Anda menguasai Kriptografi.

\vspace{1cm}
\begin{flushright}
Penyusun\\
\mbox{[Nama Dosen]}
\end{flushright}

\cleardoublepage

% ============================================================
% Cara Menggunakan Buku Ini
% ============================================================
\chapter*{Cara Menggunakan Buku Ini}
\addcontentsline{toc}{chapter}{Cara Menggunakan Buku Ini}

Buku ajar ini dirancang dengan pendekatan OBE untuk memaksimalkan pencapaian pembelajaran Anda. Berikut panduan penggunaan buku ini:

\section*{Struktur Buku}

\textbf{Bab I: Pendahuluan dan Orientasi}\\
Memperkenalkan tujuan buku, keterkaitan dengan RPS, dan konteks kurikulum OBE.

\textbf{Bab II: Landasan Teori}\\
Menyajikan fondasi teoretis kriptografi yang menjadi basis pembelajaran seluruh bab berikutnya.

\textbf{Bab III--IX: Unit Materi Inti}\\
Setiap bab mencakup satu topik utama kriptografi dengan struktur lengkap: Sub-CPMK, materi, aktivitas, latihan, asesmen, dan checklist. Implementasi menggunakan bahasa C.

\textbf{Bab X: Evaluasi dan Integrasi}\\
Berisi asesmen komprehensif dan panduan refleksi untuk mengukur pencapaian kompetensi secara menyeluruh.

\textbf{Lampiran}\\
Menyediakan rubrik penilaian, glosarium istilah kriptografi, dan referensi tambahan.

\section*{Komponen dalam Setiap Bab}

\begin{enumerate}
  \item \textbf{Sub-CPMK}: Baca dengan seksama untuk memahami kompetensi yang harus dicapai
  \item \textbf{Materi Pokok}: Pelajari dengan cermat, jalankan semua contoh kode C
  \item \textbf{Aktivitas Pembelajaran}: Lakukan secara mandiri atau berkelompok
  \item \textbf{Latihan}: Kerjakan untuk menguji pemahaman Anda
  \item \textbf{Asesmen}: Gunakan untuk mengukur pencapaian Sub-CPMK
  \item \textbf{Checklist}: Centang setelah yakin menguasai setiap indikator
\end{enumerate}

\section*{Tips Belajar Efektif}

\begin{itemize}
  \item Jangan hanya membaca, praktikkan semua contoh kode C
  \item Gunakan compiler C (GCC, Clang) atau IDE (Code::Blocks, Visual Studio Code) untuk eksperimen
  \item Kerjakan latihan sebelum melihat solusi
  \item Diskusikan konsep yang sulit dengan teman atau dosen
  \item Manfaatkan checklist untuk self-assessment berkala
  \item Kerjakan proyek implementasi kriptografi untuk mengintegrasikan konsep yang dipelajari
\end{itemize}

\cleardoublepage

% ============================================================
% Identitas Mata Kuliah
% ============================================================
\chapter*{Identitas Mata Kuliah}
\addcontentsline{toc}{chapter}{Identitas Mata Kuliah}

\begin{tabular}{ll}
  Nama Program Studi & : Matematika \\
  Nama Mata Kuliah & : Kriptografi \\
  Kode Mata Kuliah & : MAT-XXX \\
  Semester & : Ganjil/Genap (sesuai kurikulum) \\
  SKS / Bobot Kredit & : 3 SKS (2 Teori, 1 Praktikum) \\
  Jumlah Pertemuan & : 16 pertemuan \\
  Dosen Pengampu & : \mbox{[Nama Dosen]} \\
  Prasyarat & : Aljabar Linear, Matematika Diskrit \\
  Tanggal Penyusunan & : [Tanggal] \\
\end{tabular}

\vspace{1cm}

\section*{Capaian Pembelajaran Lulusan (CPL)}

CPL yang dibebankan pada mata kuliah ini mencakup kompetensi lulusan dalam aspek pengetahuan, keterampilan, dan sikap:

\begin{enumerate}
  \item \textbf{CPL-1 (Pengetahuan):} Menguasai konsep teoretis matematika termasuk aljabar, analisis, dan matematika diskrit untuk menyelesaikan masalah.
  
  \item \textbf{CPL-2 (Pengetahuan):} Menguasai prinsip-prinsip komputasi dan pemrograman untuk implementasi solusi matematika.
  
  \item \textbf{CPL-3 (Keterampilan Umum):} Mampu menerapkan pemikiran logis, kritis, sistematis, dan inovatif dalam konteks pengembangan ilmu pengetahuan.
  
  \item \textbf{CPL-4 (Keterampilan Khusus):} Mampu merancang dan mengimplementasikan algoritma matematika menggunakan bahasa pemrograman.
  
  \item \textbf{CPL-5 (Keterampilan Khusus):} Mampu menganalisis dan menyelesaikan masalah matematika terapan.
  
  \item \textbf{CPL-6 (Sikap):} Menunjukkan sikap bertanggung jawab atas pekerjaan di bidang keahliannya secara mandiri dan mampu bekerja sama dalam tim.
\end{enumerate}

\section*{Capaian Pembelajaran Mata Kuliah (CPMK)}

Kemampuan atau kompetensi spesifik yang diharapkan mahasiswa kuasai setelah menyelesaikan mata kuliah:

\begin{enumerate}
  \item \textbf{CPMK-1:} Mahasiswa mampu menjelaskan prinsip-prinsip dasar kriptografi dan konsep keamanan informasi (confidentiality, integrity, authenticity).
  
  \item \textbf{CPMK-2:} Mahasiswa mampu menerapkan aritmetika modular, algoritma Euclid, dan teori bilangan dalam konteks kriptografi.
  
  \item \textbf{CPMK-3:} Mahasiswa mampu menganalisis dan mengimplementasikan algoritma kriptografi klasik (Caesar, Vigenère, substitusi) dalam bahasa C.
  
  \item \textbf{CPMK-4:} Mahasiswa mampu menganalisis dan mengimplementasikan algoritma kriptografi modern simetris (AES, DES) dan asimetris (RSA) dalam bahasa C.
  
  \item \textbf{CPMK-5:} Mahasiswa mampu menerapkan fungsi hash dan tanda tangan digital dalam keamanan data.
  
  \item \textbf{CPMK-6:} Mahasiswa mampu merancang dan menganalisis protokol kriptografi serta manajemen kunci.
  
  \item \textbf{CPMK-7:} Mahasiswa mampu melakukan kriptanalisis dasar terhadap cipher sederhana (frequency analysis, Kasiski examination).
  
  \item \textbf{CPMK-8:} Mahasiswa mampu mengimplementasikan algoritma kriptografi menggunakan bahasa pemrograman C.
\end{enumerate}

\section*{Matriks Keterkaitan CPL dan CPMK}

\begin{table}[htbp]
\centering
\small
\begin{tabular}{|c|c|>{\raggedright\arraybackslash}p{5.5cm}|}
  \hline
  \textbf{CPL} & \textbf{CPMK} & \textbf{Keterangan} \\
  \hline
  CPL-1 & CPMK-1, 2, 3, 4, 5, 6, 7 & Penguasaan konsep teoretis kriptografi \\
  \hline
  CPL-2 & CPMK-3, 4, 8 & Prinsip komputasi dan implementasi \\
  \hline
  CPL-3 & CPMK-1 s/d 8 & Pemikiran logis dan inovatif \\
  \hline
  CPL-4 & CPMK-3, 4, 8 & Implementasi algoritma dalam C \\
  \hline
  CPL-5 & CPMK-2, 3, 4, 5, 6, 7 & Analisis masalah matematika terapan \\
  \hline
  CPL-6 & CPMK-8 & Kerja sama dalam proyek implementasi \\
  \hline
\end{tabular}
\caption{Matriks Keterkaitan CPL dan CPMK}
\end{table}

\cleardoublepage

% ============================================================
% Daftar Isi
% ============================================================
\phantomsection
\addcontentsline{toc}{chapter}{Daftar Isi}
\tableofcontents
\cleardoublepage

\clearpage
\phantomsection
\addcontentsline{toc}{chapter}{Daftar Gambar}
\sloppy
\listoffigures
\fussy
\cleardoublepage

\phantomsection
\addcontentsline{toc}{chapter}{Daftar Tabel}
\listoftables
\cleardoublepage

\phantomsection
\addcontentsline{toc}{chapter}{Daftar Kode Program}
\lstlistoflistings
\cleardoublepage
